% Verbatim excerpts from main_final.tex in the arXiv:2508.20007v2 source archive
% (https://arxiv.org/src/2508.20007v2, CC BY 4.0). Full-file SHA-256:
% 3564d5b3cbc984520fce5f0527123e276c6698db6d69579b2bed21511f0ec63c
% Each block is headed by its 1-based line range in that file.

%% lines 104-114
\begin{abstract}
 
A subparsec supermassive black hole binary (SMBHB) at the center of the galaxy 3C~66B is a promising candidate for continuous gravitational-wave searches with pulsar timing arrays (PTAs). 
In this work, we search for such a signal in the third data release of the Parkes Pulsar Timing Array. 
Matching our priors to estimates of binary parameters from electromagnetic observations, we find a log Bayes factor $\ln \mathcal{B}= - 0.0027(7)$, highlighting that the source can be neither confirmed nor ruled out. 
We place upper limits at $95\%$ credibility on the chirp mass $\mathcal{M} < 6.90 \times 10^{8}\ M_{\odot}$, and on the characteristic strain amplitude  $\textrm{log}_{10}(h_0)< -14.44$. 
This partially rules out the parameter space suggested by electromagnetic (EM) observations of 3C~66B. 
We also independently reproduce the calculation of the chirp mass with the 3 mm flux monitor data from the unresolved core of 3C~66B.
Based on this, we outline a new methodology for constructing a joint likelihood of EM and gravitational-wave data from SMBHBs.
Finally, we suggest that targeted searches may allow firmly established SMBHB candidates to be treated as standard sirens, for complementary constraints on the Universe expansion rate. 
\end{abstract}

%% lines 146-150

One of the most promising sources of CGWs is the SMBHB candidate in the center of galaxy 3C~66B. 
By using long baseline interferometry, \citet{SudouIguchi2003} found evidence for elliptical binary orbital motion in this radio galaxy's core and inferred a GW frequency of $f_{\textrm{GW}} = 60.4 \pm 1.73$ nHz. 
Further observations by \citet{IguchiOkuda2010} found that the binary has a chirp mass $\mathcal{M}=7.9^{+3.8}_{-4.5} \times 10^{8}\ M_{\odot}$, primarily through measuring the orbital radius of the electromagnetic (EM)-emitting black hole and the semi-major axis of an SMBHB.
Galaxy 3C~66B is located at a redshift $z = 0.02126$ \citep{vandenBoschGebhardt2015}, which corresponds to a luminosity distance $D_{\text{L}} = 94.17^{+0.20}_{-0.20}$~Mpc~\citep{PlanckCollaborationAghanim2020}. 

%% lines 166-178

\begin{tabular}{lll}
\hline
 & Value & \parbox[t]{0.32\linewidth}{Reference} \\ \hline \hline
 %& Value & Reference \\ \hline \hline
$\mathcal{M}$ ($M_{\odot}$) & $7.9^{+3.8}_{-4.5} \times 10^{8}$ & \citet{IguchiOkuda2010} \\
$f_{\textrm{GW}}$ (nHz) & $60.4 \pm 1.73$ & \citet{SudouIguchi2003} \\
$D_\text{L}$ (Mpc) & $94.17 \pm 0.20$ & \citet{vandenBoschGebhardt2015} \\
\shortstack[l]{R.A. (HH:MM:\\SS.SSS)} & 02:23:11.4110 & \citet{HuntJohnson2021} \\
\shortstack[l]{Decl. (DD:MM:\\SS.SSS)} & +42:59:31.385 & \citet{HuntJohnson2021} \\
$h_0$ & $7.3^{+6.8}_{-5.8} \times 10^{-15}$ & \citet{IguchiOkuda2010} \\
\hline
\end{tabular}

%% lines 353-356
Similarly, we fix the luminosity distance based on the redshift of 3C~66B, assuming the $\Lambda$CDM cosmological model based on~\citet{PlanckCollaborationAghanim2020}.
We use an exponential-Normal prior for $\log_{10} f_\text{GW}$, with mean and variance from Table~\ref{tab:3c66b}. 
We perform this calculation with only the Earth term of the CGW.
Finding $\ln \mathcal{B} = - 0.0027(7)$, we claim no evidence for an SMBHB at 3C~66B; nevertheless, we cannot rule out the EM model. 

%% lines 370-380

The results are shown in Figure \ref{fig:MassPosteriors}. 
When only the Earth term is modeled in our analysis, we find an upper limit of $\mathcal{M} < 8.25 \times 10^{8}~M_\odot$ at 95\% credibility. 
When including the pulsar term, we find a more constraining limit, $\mathcal{M} < 6.90 \times 10^{8}~M_\odot$. 
For comparison, we show the value from the EM model from~\citet{IguchiOkuda2010}, $(7.9^{+3.8}_{-4.5})\times10^{8}$ $M_{\odot}$. 
Our limits are not below the prediction of the EM model. 
Therefore, we cannot rule it out.
However, the tightest upper limit calculated rules out 58\% of the uncertainty region of the EM-model. 
For consistency with~\cite{ng12_allsky_cw} and \cite{ng15_cw_targeted}, we also obtain a posterior by reweighting our posteriors to a uniform prior in linear space $\pi(\mathcal{M})=\mathcal{U}(10^{7.5},10^{9.5})~[M_\odot]$. 
While not matching astrophysical expectations for the SMBHB mass function, this presents the most conservative limits. 
With this prior, we report an upper limit at $95\%$ credibility to be $\mathcal{M} < 11.06 \times 10^{8} ~M_\odot$ for the Earth term and $\mathcal{M} < 10.63\ \times 10^{8} ~M_\odot$ when the pulsar term is included. 

%% lines 412-414
With only the Earth term, we find an upper limit at 95\% credibility to be $\log_{10} h_0 < -14.12$. 
When including the pulsar term, we find a more constraining limit, $\log_{10} h_0 < -14.44$. 
Using the more conservative prior, $\pi(h_0)=\mathcal{U}(10^{-17.0},10^{-12.5})$, we find an upper limit for the Earth term of $\log_{10} h_0 < -13.86$ and an upper limit of $\log_{10} h_0 < -14.20$ when including the pulsar term. As in Section~\ref{sec:results:mc}, the range of $h_0$ of the EM model from~\citet{IguchiOkuda2010} is partially ruled out in all cases.

%% lines 476-488
\subsection{\label{sec:results:frequentist} Frequentist Search}

To confirm the nondetection of a CGW, we perform a frequentist targeted search based on the $\mathcal{F}$ statistic~\citep{EllisSiemens2012}, a frequentist measure of evidence for a CGW in our data. 
In particular, we chose to calculate the $\mathcal{F}_\text{e}$ statistic, which is based on only modeling the Earth term of the CGW signal. 
This statistic is a function of CGW frequency and sky position, which we set to the 3C~66B EM model values.
We find a value of $\mathcal{F}_\text{e}=4.90$, corresponding to the false alarm probability $\text{FAP}=0.50$. 
\citet{EllisSiemens2012} suggest a threshold for FAP to be $10^{-4}$ to claim a detection.
The trials factor for our search is one.
The background distribution of the $\mathcal{F}_\text{e}$ statistic is the chi-squared distribution with four degrees of freedom and is displayed in Figure \ref{fig:FeStat_fulltarget} along with the value of the $2\mathcal{F}_\text{e}$ statistic. 
Based on this null distribution, we find a $p\text{-value} = 0.04$ which corresponds to about $2.1 \sigma$.
Thus, the $\mathcal{F}_\text{e}$ statistic value we find is consistent with noise. 

 \begin{figure}[h]

%% lines 557-570
\section{\label{sec:conclusion} Conclusion}

Using PPTA DR3, we conduct a targeted multimessenger search for an SMBHB at the center of 3C~66B. 
We search for a CGW signal by performing model selection, where we compare a model with a CGW signal to the null hypothesis. 
With priors carefully matched to the EM model, we find a log Bayes factor of $\ln \mathcal{B} = -0.0027(7)$. 
This suggests we can neither confirm nor rule out the EM model. 
We also find a frequentist statistic, $\mathcal{F}_e = 5.04$, with an associated $\text{FAP}=0.50$, which confirms nondetection of the signal. 
Based on the background estimation, we find that the value of $\mathcal{F}_e = 5.04$ corresponds to $p=0.04$, about $2.1\sigma$. 

Following nondetection, we place upper limits on the mass as $\mathcal{M} < 6.90 \times 10^8~M_{\odot}$  and strain amplitude $\log_{10}(h_0) < -14.44$, at 95\% credibility. 

For both the mass and the strain, we rule out over 50\% of the parameter space of the EM model. 
Using conservative priors, we still rule out a portion of the EM model, finding a $95\%$ upper limit on the chirp mass of $\mathcal{M} < 9.55 \times 10^8~M_{\odot}$ and for the strain amplitude we find a value of $h_0 < -14.20 $. 
Finally, we demonstrate a new methodology for a joint-likelihood approach for the analysis of 3C~66B as well as a proof of principle, showing that it may be possible to use firmly established SMBHB candidates in targeted searches with PTAs as standard sirens, complementary to high-precision cosmology.

%% lines 706-762
\section{\label{sec:appendix:pdist}Pulsar Distances}

Pulsar distances are required in the calculation of the expected signal when modeling the pulsar term of the CGW signal. 
Pulsar distances are measured using very long baseline interferometry (VLBI) astrometry, optical astrometry, timing, and the \citet{Shklovskii1970} effect. 
A list of the priors used for each pulsar, as well as the technique used to find this distance, are listed in Table~\ref{tab:pulsar_distances}. 
We impose a Gaussian prior on pulsar distance in accordance with the estimates found in the literature. 
When no such estimates are available for a given pulsar, we impose a sufficiently conservative prior $\mathcal{N}(1000,200)~\text{[pc]}$. 
Distance estimates based on pulsar dispersion measures are not used due to the large uncertainties in the electron density model~\citep{YaoManchester2017}. 

\begin{table*}[]
%\scriptsize
\renewcommand{\arraystretch}{1.3}
\centering
\caption{Table Describing the Pulsar Distances, the Methods Used to Obtain Them, and the Source from Which They Were Obtained.}
\begin{tabular}{cccc}
\hline
PSR & Distance [pc]                     & Measurement             & Reference \\ \hline
 & (pc)                     &              & \\ \hline \hline
J0030+0451  & $329^{+6}_{-5}$               & VLBI astrometry                   & \citet{DingDeller2023}       \\
J0125-2327  & $1000 \pm 200$                 & ...                                 & \citet{EllisVanHaasteren2017}       \\
J0437-4715  & $156.96\pm0.11$               & Shklovskii effect                 & \citet{Shklovskii1970, ReardonBailes2024}       \\
J0613-0200  & $125\pm54.69$                 & Timing                            & \citet{VerbiestBailes2009}       \\
J0614-3329  & $1000 \pm 200$                 & ...                                 & \citet{EllisVanHaasteren2017}       \\
J0711-6830  & $1000 \pm 200$                 & ...                                 & \citet{EllisVanHaasteren2017}       \\
J0900-3144  & $1000 \pm 200$                 & ...                                 & \citet{EllisVanHaasteren2017}       \\
J1017-7156  & $256.41\pm78.90$              & Timing                            & \citet{NgBailes2014}       \\
J1022+1001  & $720.98^{+21.31}_{-14.55}$    & VLBI astrometry                   & \citet{DellerGoss2019}       \\
J1024-0719  & $1072^{+67}_{-49}$            & Optical astrometry                & \citet{MoranMingarelli2023}       \\
J1045-4509  & $303.03\pm 174.47$            & Timing                            & \citet{VerbiestBailes2009}       \\
J1125-6014  & $1000 \pm 200$                 & ...                                 & \citet{EllisVanHaasteren2017}       \\
J1446-4701  & $1000 \pm 200$                 & ...                                 & \citet{EllisVanHaasteren2017}       \\
J1545-4550  & $1000 \pm 200$                 & ...                                 & \citet{EllisVanHaasteren2017}       \\
J1600-3053  & $5000 \pm 3750$               & Timing                            & \citet{VerbiestBailes2009}       \\
J1603-7202  & $1000 \pm 200$                 & ...                                 & \citet{EllisVanHaasteren2017}       \\
J1643-1224  & $4763.36^{+99.06}_{-104.89}$  & VLBI astrometry                   & \citet{DingDeller2023}       \\
J1713+0747  & $1052.63^{+66.48}_{-55.40}$   & VLBI astrometry                   & \citet{ChatterjeeBrisken2009}       \\
J1730-2304  & $510^{+30}_{-30}$             & VLBI astrometry                   & \citet{DingDeller2023}       \\
J1744-1134  & $416.67 \pm 17.36$            & Timing                            & \citet{VerbiestBailes2009}       \\
J1824-2452A & $10000 \pm 50000$             & VLBI astrometry                   & \citet{DingDeller2023}       \\
J1832-0836  & $1000 \pm 200$                 & ...                                 & \citet{EllisVanHaasteren2017}       \\
J1857+0943  & $909.09 \pm 165.29$           & VLBI astrometry                   & \citet{VerbiestBailes2009}       \\
J1902-5105  & $1000 \pm 200$                 & ...                                 & \citet{EllisVanHaasteren2017}       \\
J1909-3744  & $1265.82 \pm 32.05$           & Timing                            & \citet{VerbiestBailes2009}       \\
J1933-6211  & $1000 \pm 200$                 & ...                                 & \citet{EllisVanHaasteren2017}       \\
J1939+2134  & $2604.17^{+325.52}_{-311.96}$ & VLBI astrometry                   & \citet{DingDeller2023}       \\
J2124-3358  & $322.58 \pm 57.23$            & Timing                            & \citet{VerbiestBailes2009}       \\
J2129-5721  & $2358.49 \pm 489.50$          & Optical astrometry           & \citet{JenningsKaplan2018}       \\
J2145-0750  & $623.83^{+24.51}_{-3.50}$     & VLBI astrometry                   & \citet{DellerGoss2019}       \\
J2241-5236  & $1000 \pm 200$                 & ...                                 & \citet{EllisVanHaasteren2017}   \\ \hline   
\end{tabular}
\vspace{2mm}
\begin{minipage}{\textwidth}
\centering
\footnotesize
Note. If no distance was available, the method is listed as ``...".
\end{minipage}
\label{tab:pulsar_distances}
